\section{Introduction}
\label{sec:intro}

A tool-calling agent often answers a request with several calls at once: the weather in two cities, a loan check at two banks. The calls of such a turn form a joint action, and they must fit together: each call needs its own entity, and no call may repeat another. Diffusion language models (dLLMs) look like natural decision cores for such agents, since they commit many tokens per denoising step and could write all the calls of a turn in a few steps. This raises a coordination worry. Tokens committed in the same step are chosen without seeing each other, like the moves of players who act at the same time and cannot talk, so the calls of a parallel turn might collide or bind each other's arguments. dLLM agents do fail at tool calling, and the failures have been blamed on diffusion noise and parallel decoding~\cite{lu2026bitter}, but the decoding protocol has not been separated from the agent's other design choices.

One such choice is easy to overlook. To keep every call well formed, a dLLM agent can write the syntax of its turn in advance as a skeleton and leave only the argument values masked. Each value then gets a slot with a fixed number of masks, set before the value is known. In deployment the slot lengths come from a fixed budget, a per-argument maximum, or an estimate, and all three are often wrong: a budget leaves surplus masks, a per-argument maximum gives one call a slot that fits a sibling's value, and an estimate does both. Figure~\ref{fig:example} shows the effect when the agent commits one token per step, so that no two decisions are ever made at once: with one extra mask per slot, HSBC becomes ``HSBC Bank'' and a \$500{,}000 loan becomes one of \$5{,}000{,}000; with the slot lengths of the two banks swapped, the banks change calls.

We treat a dLLM agent's parallel turn as simultaneous decisions by its slots and ask what these decisions coordinate on. A slot can tell its own value from its siblings' only through cues that all slots share: the order in which the request mentions its entities (a focal point), the position of the slot (a role label), and the length of the slot (a signal). From these cues we derive predictions and test them on the parallel calls of BFCL~\cite{patil2025bfcl} with two dLLMs and an AR agent, varying the decoding protocol and the slot lengths independently. To check the analogy, we give symmetric requests to teams of LLM agents that act at the same time or take turns.

We find three things. (1) When the request orders its entities, the slots follow the mention order at any degree of parallelism, so committing 16 tokens per step instead of one costs \dream only 5.3 points of set accuracy. (2) When nothing tells the calls apart, slots committed in the same step collide, in every open choose-N request against none with one token per step. Teams of LLM agents collide in the same way when they act at the same time and never when they take turns; unlike the slots, they are not kept apart by role numbers. (3) The costliest failures come from the interface. The models read the number of masks in a slot as a promise about the length of its value: one extra mask per slot lowers \dream's set accuracy from 89.8\% to 15.5\%, and swapped slot lengths make 79.4\% of the affected slots take the sibling value that fits, with one token per step. With slot lengths that the model estimates itself, one token per step produces about four times as many cross-call errors as 16 tokens per step with exact lengths (7.5\% against 2.0\% of \dream's requests): errors that an evaluator would blame on parallel decoding.

Our contributions are (i) a coordination analysis of a dLLM agent's parallel tool calls, with predictions tested on the agent and on teams of LLM agents (Sections~\ref{sec:setup}--\ref{sec:symmetric}); (ii) evidence that masks act as length promises, with near misses the most harmful (Section~\ref{sec:length}); (iii) evidence that wrong slot lengths produce the cross-call errors attributed to parallel decoding, and a protocol for attributing errors correctly (Section~\ref{sec:masquerade}); and (iv) design guidance for dLLM agents (Section~\ref{sec:design}). Length biases of dLLMs are known for single spans~\cite{wu2026dreamon,han2026dia,liu2026cal}; we show what they do to an agent's joint actions.
